\documentclass[11pt]{article}
\usepackage[margin=1in]{geometry}
\usepackage[T1]{fontenc}
\usepackage[utf8]{inputenc}
\usepackage{lmodern}
\usepackage{enumitem}
\usepackage{xcolor}
\usepackage{hyperref}
\usepackage{setspace}
\usepackage{titlesec}
\usepackage{parskip}

\definecolor{NYUPurple}{HTML}{57068C}
\hypersetup{colorlinks=true, linkcolor=NYUPurple, urlcolor=NYUPurple}
\titleformat{\section}{\large\bfseries\color{NYUPurple}}{}{0em}{}
\titleformat{\subsection}{\normalsize\bfseries}{}{0em}{}

\title{Practice Transcript for Presentation \\ \large Capacity Constraints, Case-load Increases and Incarceration Rates}
\author{Prepared for rehearsal}
\date{April 2026}

\begin{document}
\maketitle
\onehalfspacing

\section*{How to use this document}
This is not meant to be read word-for-word. It is a practice script. The goal is to help you rehearse the flow, tighten transitions, and remember the main interpretation and caveats. I wrote it in a way that sounds formal but still natural when spoken.

I would practice it in three passes:
\begin{itemize}[leftmargin=1.5em]
    \item First pass: read it almost verbatim to internalize the structure.
    \item Second pass: shorten each paragraph into your own words while keeping the transitions.
    \item Third pass: practice only from the bolded takeaways and transition lines.
\end{itemize}

\section{Opening}
\textbf{Possible opening:}

Thank you very much for being here. This paper studies how criminal justice actors respond when institutional constraints are relaxed. In particular, I focus on prison capacity. The core question is whether exogenous increases in prison space change judicial and policing behavior, and therefore affect incarceration rates.

The broader motivation is that criminal justice decisions are often treated as if they were driven only by crime and legal criteria. But in practice, they may also reflect institutional constraints. If that is true, then changing capacity may change outcomes even if underlying crime does not.

\textbf{Transition:} So I start by motivating why this is an important setting to study.

\section{Slide 1: Motivation}
Institutional actors respond to constraints. We see this in many settings: bureaucrats respond to targets, teachers respond to testing incentives, and firms respond to regulation and taxes. Criminal justice should not be different.

In this setting, there are several possible constraints. Some are fiscal: incarceration is costly. Some are electoral: in some contexts, criminal justice actors respond to political incentives. Some are administrative: rising caseloads leave less time to process cases. And some are capacity-related: when prisons are full, incarceration becomes harder to use.

The key idea of the paper is simple. If criminal justice behavior changes when these constraints bind, then relaxing one of them may also change outcomes. So my research question is whether exogenous increases in prison capacity affect judicial and policing behavior.

\textbf{What to emphasize verbally:}
\begin{itemize}[leftmargin=1.5em]
    \item You do not need to spend much time on all four constraints.
    \item The point is just to tell the audience that prison capacity is one example of a broader class of institutional constraints.
    \item End clearly on the research question.
\end{itemize}

\textbf{Transition:} Once we think of capacity in this way, the main challenge becomes identification.

\section{Slide 2: Strategy}
Causal evidence is scarce for two reasons.

First, criminal justice outcomes are shaped by many actors. Judges matter, but so do prosecutors, police, and prison authorities. So when incarceration changes, it is not obvious which margin is driving the result.

Second, prison construction is usually endogenous. Governments tend to expand capacity where overcrowding is already severe or where incarceration demand is already rising. So a simple comparison between places with and without new prisons would not have a causal interpretation.

To address this, I exploit the staggered construction of federal prisons across Mexico as plausibly exogenous shocks to local prison capacity. The empirical strategy relies on three institutional features. Judges must assign inmates to the nearest available facility. New federal prisons were built to relocate federal inmates, not to solve local capacity problems. And state prisons also house federal prisoners. Together, these features imply that federal prison expansion can relieve pressure on local state prisons.

\textbf{Short version if you need to move fast:}
The strategy is to use staggered federal prison construction as an exogenous shock to local capacity, leveraging the fact that judges assign inmates to nearby facilities and that state prisons also housed federal prisoners.

\textbf{Transition:} Before showing the results, I give some institutional context for why these shocks matter locally.

\section{Slide 3: Capacity vs. population figure}
This figure shows the evolution of prison capacity relative to the prison population. The main takeaway is that overcrowding is a real and persistent feature of the Mexican prison system during this period.

What I want the audience to take from this figure is not just that prisons are crowded, but that capacity is a meaningful institutional constraint. When the system is operating under pressure, expanding capacity may affect behavior throughout the criminal justice chain.

\textbf{Transition:} The next three slides explain why federal prison construction can affect local outcomes.

\section{Slide 4: Inmate assignment}
The first key institutional feature is inmate assignment. Judges are required to assign inmates to the nearest available facility relative to the municipality of residence. This matters because it makes distance to a new prison substantively meaningful. If a new prison opens nearby, it changes the relevant local incarceration margin.

\textbf{Transition:} The second institutional feature is the federal expansion itself.

\section{Slide 5: Federal prison expansion}
In 2009, under President Felipe Calderón, a federal mandate required individuals charged with federal offenses to be transferred to federal prisons. This expanded federal capacity and shifted inmates out of state facilities.

The important point for my design is that these new federal prisons were not built to solve local state-level overcrowding. Their purpose was to relocate federal inmates. But because state prisons also housed federal prisoners, this expansion still relieved pressure on local facilities.

\textbf{Tip for delivery:}
Pause on the map and say something like: ``So the treatment is local exposure to these new federal prisons, not simply federal prison construction in the abstract.''

\textbf{Transition:} The third feature is that federal and state systems were interconnected.

\section{Slide 6: Housing of federal prisoners in state facilities}
This slide shows that state prisons housed a non-trivial share of federal prisoners. This is what connects federal expansion to local state prison congestion.

So the mechanism is not that local judges suddenly gain direct access to new federal prisons for ordinary state cases. The mechanism is that federal prisoners are shifted out of state facilities, freeing up capacity locally.

\textbf{Transition:} With that institutional setup in mind, here is a preview of the main findings.

\section{Slide 7: Preview of findings}
The first result is a validation result. I show that federal prison construction decreased overcrowding in state prisons. That is important because it confirms that the treatment is actually shifting local capacity.

The main result is that federal prison construction increases pretrial processing and the number of individuals sentenced.

More specifically, municipalities closer to new federal prisons experience higher local incarceration rates, a greater likelihood of pretrial detention, and no change in monetary fines.

The broad interpretation is that the results are consistent with changes in judicial behavior driven by capacity constraints, rather than solely by increased caseloads or resource constraints. At the same time, I am careful not to overstate this because those alternative channels cannot be fully ruled out.

\textbf{Good one-sentence summary:}
Relaxing prison capacity appears to increase the use of incarceration as a margin of punishment.

\textbf{Transition:} The next two slides place these findings in the broader theoretical literature and clarify how I interpret them.

\section{Slide 8: Theory and context}
The literature suggests three broad channels through which criminal justice outcomes may respond to institutional constraints.

The first is salience. Incarceration may become more likely when prison is a more salient or central punishment.

The second is fiscal or capacity constraints. When incarceration is costly or prison space is limited, actors may substitute away from prison.

The third is time constraints. Higher caseloads may change how judges and prosecutors process cases.

In my setting, the most relevant of these channels is capacity. Nearby federal prison construction relaxes local prison constraints. By contrast, the electoral and local fiscal incentives that appear in some US studies are less likely to operate in the same way here.

\textbf{What to emphasize:}
This slide is not repeating the motivation. It is narrowing the relevant theories for your context.

\textbf{Transition:} Given the findings, the question becomes how to interpret the pattern and what caveats remain.

\section{Slide 9: Interpretation and caveats}
The findings show more pretrial processing, more pretrial detention, more prison sentences, and no change in monetary fines. That pattern suggests that relaxing capacity increases the use of incarceration.

This is more consistent with a capacity-based explanation than with a purely administrative response. If this were only an administrative story, I might expect more case movement overall without such a clear shift toward incarceration specifically.

At the same time, an important caveat remains. Capacity expansions may also increase policing and case inflows. So part of the effect may operate through higher caseload pressure rather than judicial behavior alone.

My interpretation, therefore, is not that I fully isolate one actor. Rather, I interpret the results as evidence that relaxing capacity constraints changes criminal justice outcomes, while the precise margin remains harder to fully disentangle.

\textbf{Good line to memorize:}
The evidence points clearly to an effect of capacity, but less clearly to the exact actor-level mechanism.

\section{Slide 10: Data}
The analysis uses two administrative datasets covering all formally processed individuals in Mexico at the local level.

The first is a pretrial dataset that captures early judicial decisions, especially detention versus release. The second is a criminal courts dataset that captures final outcomes such as prison sentences, monetary fines, and acquittals.

A major advantage is that both datasets include the municipality of residence. That allows me to aggregate the data to the municipality-by-bimonth level and match them to geolocated data on prison capacity and prison population.

\textbf{Transition:} With those data in hand, I estimate a staggered difference-in-differences design.

\section{Slide 11: Empirical strategy}
I estimate a generalized difference-in-differences model with staggered adoption. The treatment variable turns on when the closest federal prison to a municipality is within 300 kilometers.

Intuitively, the estimate compares how outcomes change in municipalities that become exposed to a nearby federal prison before and after the opening, relative to municipalities that are not exposed over the same period.

I use log transformations for count outcomes to deal with skewness and outliers, and I show that the results are robust to alternative transformations and thresholds.

\textbf{Tip:}
Do not spend too much time on the equation unless someone in the audience is very methods-focused. The key is the intuition.

\textbf{Transition:} I then move to the results.

\section{Slide 12: Overcrowding event study}
This is the first-stage validation result. After nearby federal prison openings, relative overcrowding in local state prisons declines.

This matters because it shows that the treatment is doing what the design requires it to do: it is relaxing local capacity constraints.

\textbf{Good line:}
Before interpreting sentencing effects, I first verify that the treatment actually changes local prison congestion.

\section{Slide 13: Processing table}
The first set of outcome results concerns processing. I find increases in the total number of processed cases, increases in pretrial detention, and increases in releases in levels.

But in proportional terms, what stands out most is the increase in pretrial detention relative to processed cases. So the treatment appears to affect both the extensive margin of case processing and the conditional use of detention.

\textbf{How to phrase carefully:}
This is consistent with greater reliance on detention when capacity is relaxed, though some of the increase in processing may also reflect higher case inflows.

\section{Slide 14: Processing event studies}
The event studies are useful for two reasons. First, they let me visually assess pre-trends. Second, they show the timing of the response.

The pattern is that the effect emerges after treatment and not before, which supports the identification strategy. The main message is that the dynamic response is consistent with an effect of nearby prison construction on pretrial processing behavior.

\section{Slide 15: Sentencing table}
Turning to final sentencing outcomes, I find an increase in the total number sentenced and in guilty verdicts that lead to prison, but essentially no change in guilty verdicts that lead to monetary fines.

That distinction is important. It suggests that the effect is not simply a general increase in punishment, but a more specific increase in incarceration as the chosen form of punishment.

\textbf{Good line:}
The margin that moves is prison, not punishment in general.

\section{Slide 16: Sentencing event studies}
The event-study results reinforce the table evidence. The post-treatment pattern is concentrated in prison-related outcomes rather than monetary fines.

This strengthens the interpretation that the policy shock changes the use of incarceration specifically.

\section{Conclusion slide}
To conclude, I provide evidence that exogenous increases in prison capacity affect criminal justice outcomes. Federal prison construction relieves overcrowding in local state prisons and increases pretrial processing, pretrial detention, and prison sentences.

The broader implication is that incarceration rates may be shaped not only by underlying crime and legal criteria, but also by institutional capacity. In that sense, prison construction may do more than absorb existing prisoners; it may change the behavior of the criminal justice system itself.

Thank you very much, and I look forward to your comments.

\section{Possible questions and short answers}
\subsection*{1. How do you know this is not just more crime?}
My design compares municipalities before and after nearby federal prison openings relative to municipalities not exposed over the same period. The identifying variation comes from staggered prison construction, not from raw differences in crime levels. Also, the fact that fines do not move while prison outcomes do suggests a specific incarceration mechanism rather than a generic crime shock.

\subsection*{2. Is this really judges, or could it be police?}
That is an important caveat. I interpret the results as evidence that relaxing capacity changes criminal justice outcomes, but I cannot fully separate judicial responses from policing responses. The data suggest both margins may matter.

\subsection*{3. Why Mexico?}
Mexico is useful because it provides a setting where federal prison construction generates plausibly exogenous local capacity shocks, and where some of the electoral and local fiscal mechanisms emphasized in the US literature are less central.

\subsection*{4. Why should we care substantively?}
Because if incarceration responds to capacity, then incarceration rates are not pinned down only by crime or legal merit. Institutional constraints can shape who gets detained and who gets sentenced to prison.

\subsection*{5. What is the main caveat?}
The main caveat is that capacity expansions may also increase policing and case inflows, so part of the estimated effect may reflect higher caseload pressure rather than judicial behavior alone.

\section{Final advice for rehearsal}
\begin{itemize}[leftmargin=1.5em]
    \item Keep returning to the same core sentence: \textbf{relaxing capacity increases the use of incarceration.}
    \item Do not over-explain every slide. Your talk is strongest when each slide has one clear takeaway.
    \item On theory, emphasize that you are narrowing the relevant channels, not reviewing the entire literature.
    \item On caveats, be transparent but do not undersell the result. The effect of capacity is clear even if the exact actor-level mechanism is not fully isolated.
    \item End with the broad implication: \textbf{capacity does not only accommodate incarceration; it may also generate more of it.}
\end{itemize}

\end{document}
